\documentclass[11pt]{article}
\usepackage[a4paper,margin=26mm]{geometry}
\usepackage{amsmath,amssymb,amsthm,hyperref}
\newtheorem{theorem}{Theorem}
\newtheorem{lemma}{Lemma}
\newcommand{\E}{\mathbb E}
\title{Sharp relative balance at the endpoint}
\author{Complete proof; independently reviewed}
\date{30 September 2026}
\begin{document}\maketitle
Assume positive row weights $w$ and nondecreasing columns
$p_i\in[0,1]$, $i\in[m]$, satisfy
$\E_w\prod_{i\in S}p_i=1/(|S|+1)$ for every subset $S$.
Put $u_i=1-p_i$ and $s=\sum_i u_i$.

\begin{theorem}\label{main}
There is an absolute constant $C$ such that every row satisfies
\[
 \max_i u_i\le\frac C m(s+m^{-1}).
\]
The order $m^{-1}$ of the additive term cannot be replaced by
$o(m^{-1})$ uniformly in this tensor class; see
\path{endpoint_profile_sharpness.tex}.
\end{theorem}

This is a consequence of the reviewed quadratic-bound machinery,
not an input to its proof. We use its notation
$\ell=\lfloor m/4\rfloor$, $\alpha=\ell/m$, $x=\alpha s$,
$\bar A=\E_B\prod_{j\in B}p_j$, and the measures
$\mu=(\ell+1)w\bar A$, $\nu=mw$. The scalar target is
$b_0(z)=(\ell+1)\ell^{-1}(1-z/\ell)^\ell$ on $[0,\ell]$.
Set $\theta=1/16$ and
\[
 T_d(y)=\sum_{j<d}(1-j/d)L_j(y),\qquad
 U_d(z)=d^{-2}T_d(\theta z)^2.
\]
All constants below may depend on the fixed absolute constants in
already established estimates. Only two further iterations are used.

\begin{lemma}[A sharp endpoint distribution bound]\label{sharpCDF}
For every such tensor,
\[
 \nu([0,t])\le C(t+m^{-1})\qquad(t\ge0).
\]
\end{lemma}
\begin{proof}
The proof of the $U_d$ comparison in
\path{size_bounds/quadratic_per_die_lower_bound.tex} is uniform for
$m^{1/5}\le d\le\kappa_0m$, with $\kappa_0>0$ sufficiently small.
Indeed its small, middle, and large regions give, with
$\delta=m^{-1/3}$ and $h=m^{-3/4}$,
\[
 |\mu(U_d)-\beta_0(U_d)|
 \le C\left\{\frac{\log d}{dm}+\frac hm+
                   \frac\delta m+\frac{\delta hd}m\right\}
       +O(m^{-3}).                                      \tag{1}
\]
To check the uniformity, the small-region contribution is
$C(d/m)(1/d+\delta)(1/d+h)$, bounded by the displayed terms.
The good middle-region envelope is
$C(dm)^{-1}e^{-cx}(x^{-1}+\delta x^{-2})$.
The bad-subset exponent is $-cmx/(x+\delta)$ for $x>1/d$;
it dominates every polynomial factor throughout this degree range.
The large-region tail has a polynomial factor times $e^{-cd}$.
These are exactly the estimates used in the cited proof, with the
same choice of a small upper bound on $d/m$.

Multiplying (1) by $d$ gives $o(1)$ uniformly: the largest term is
$\delta h d^2/m\le\kappa_0^2m^{-1/12}$.
Thus $\mu(U_d)\le C/d$, since $\beta_0(U_d)\le C/d$.
Also $U_d\ge c$ on $[0,c_0/d]$, and $\nu\le C\mu$ there because
$\bar A\ge1-s\ge1/2$. Choose $d$ comparable to
$\min(\kappa_0m,1/t)$ whenever $t\lesssim m^{-1/5}$.
The previous bound $\nu([0,t])\le C(t+m^{-3/4})$ handles larger $t$.
This proves the assertion, including a possible atom at zero.
\end{proof}

We next state a quantitative improvement mechanism. Suppose that,
uniformly for tensors of every sufficiently large dimension $m$,
\begin{equation}\label{hyp}
 \max_i u_i\le C_*(s+\delta_m)/m,\qquad
 \delta_m=m^{-a},\quad 1/3\le a\le1.
\end{equation}
The constant $C_*$ is fixed. Deleting one column gives the same
hypothesis with dimension $M=m-1$, and $\delta_M\asymp\delta_m$.
For a fixed deleted index $i$, let
$s_-=\sum_{j\ne i}u_j$, $\ell=\lfloor M/4\rfloor$,
$x=(\ell/M)s_-$ and let $\bar A$ average the success products
over $\ell$-subsets of these $M$ columns. Define
\[
 \mu_0=(\ell+1)w\bar A,\qquad
 \mu_1=(\ell+1)(\ell+2)w u_i\bar A,
 \qquad \nu=Mw.
\]
These are probability measures except that $\nu$ has mass $M$.
Their scalar targets satisfy
$b_1(z)=((\ell+2)/\ell)z b_0(z)$, and
$d\mu_1=(\ell+2)u_i\,d\mu_0$.
Absorbing the $u_i$ term in \eqref{hyp} gives
\begin{equation}\label{mark}
 M u_i\le C(s_-+\delta),\qquad\delta=M^{-a}.
\end{equation}

\begin{lemma}[Marked taper comparison]\label{marked}
There is a sufficiently small fixed $\kappa_0>0$ such that, uniformly
for $M^{1/5}\le d\le\kappa_0M$,
\begin{align}
 |\mu_0(U_d)-\beta_0(U_d)|
 &\le C\left\{\frac{\log d}{dM}+\frac1{M^2}
                  +\frac\delta M+\frac{\delta d}{M^2}\right\}
        +O(M^{-3}),                                  \tag{2}\\
 |\mu_1(U_d)-\beta_1(U_d)|
 &\le C\left\{\frac{1+\delta\log d}{dM}
                         +\frac{\delta^2}M\right\}
        +O(M^{-3}).                                  \tag{3}
\end{align}
In particular $\mu_0(U_d)\ge c/d$ uniformly in this range for large
$M$, and
\[
 \mu_1(U_d)\le C/d^2
 \quad\text{if additionally }d\le \sqrt M/\delta.
                                                               \tag{4}
 \]
\end{lemma}
\begin{proof}
The exact marked polarization identity is proved in
\path{rational_designs/asymptotic/quantitative_relative_endpoint_balance.tex}.
The mark is outside all success and polarized subsets, so the same
subset moments and Taylor estimates apply. The proof of (2) is the
unmarked proof recalled in Lemma~\ref{sharpCDF}, now using
$h=M^{-1}$ from that lemma and the bound \eqref{hyp}.

For (3), the extra normalized marked factor is at most $C(x+\delta)$
by \eqref{mark}. In the small region $x\le1/d$, the integrated
comparison error is at most
\[
 C\frac dM(1/d+\delta)^2\nu([0,1/d])
 \le\frac CM(1/d+\delta)^2,
\]
because $d\le\kappa_0M$. In the good middle region the unmarked
envelope therefore becomes
\[
 \frac C{dM}e^{-cx}
          (1+2\delta/x+\delta^2/x^2).
\]
The sharp CDF bound gives
\[
 \int_{1/d}^\infty e^{-cx}\,d\nu\le C,\qquad
 \int_{1/d}^\infty e^{-cx}x^{-1}\,d\nu
       \le C(\log d+d/M),
\]
and
\[
 \int_{1/d}^\infty e^{-cx}x^{-2}\,d\nu
       \le C(d+d^2/M).
\]
Substitution proves (3). The bad-subset and large-region errors
remain $O(M^{-3})$: the extra marked factor is bounded by a
polynomial in $M$, while their bounds are superpolynomially small.
For the weighted Taylor remainder the mark is constant on the
subsets of each row, so it introduces precisely the same factor.

The unmarked scalar integral is at least $c/d$, by near-zero
positivity of $U_d$ and $b_0$. The error in (2) is $o(1/d)$ uniformly,
because $\delta\le M^{-1/3}$ and $d\le\kappa_0M$.
For the marked scalar integral, use $b_1(z)\le Cz e^{-\theta z}$ and
\[
 \int_0^\infty y e^{-y}T_d(y)^2\,dy=\frac{d+1}{2d}
\]
to get $\beta_1(U_d)\le C/d^2$.
When $d\le\sqrt M/\delta$, each term in (3) is $O(1/d^2)$;
in particular $\delta\log d=o(1)$ and $d\le\kappa_0M$.
This proves (4).
\end{proof}

\begin{lemma}[Relative-balance improvement]\label{step}
Hypothesis \eqref{hyp} implies
\[
 \max_i u_i\le\frac C m
          \left(s+m^{-1}+m^{-a-1/2}\right).
\]
\end{lemma}
\begin{proof}
We first obtain an interval with controlled density ratio. There are
fixed $0<a_0<b_0$ and a sufficiently small $\kappa>0$ such that
\begin{equation}\label{interval}
 \mu_0(a_0/d<x<b_0/d)\ge c/d,\qquad
 \mu_1(x\le b_0/d)\le C/d^2
\end{equation}
uniformly for
\[
 M^{1/4}\le d\le D:=\min(\kappa M,\sqrt M/\delta).
\]
To prove the first bound, start with $\mu_0(U_d)\ge c/d$.
For $x\le a_0/d$, the ordinary Gaussian Laguerre bound gives
$U_d(x)\le C$, so this contribution is at most
$C(a_0/d+1/M)\le C(a_0+\kappa)/d$ by Lemma~\ref{sharpCDF}.
For $x\ge b_0/d$, the all-orders estimate at order zero gives
$U_d(x)\le C e^{\theta x}/(d^2x^2)$. Since
$d\mu_0\le C e^{-x}d\nu$, its contribution is at most
\[
 \frac C{d^2}\int_{b_0/d}^\infty e^{-cx}x^{-2}\,d\nu
 \le C\left(\frac1{b_0d}+\frac1{Mb_0^2}\right).
\]
Choose $b_0$ large, then $a_0$ and $\kappa$ small. A fixed positive
fraction of $\mu_0(U_d)$ remains on the indicated interval, where
$U_d\le C$. This proves its probability lower bound.

For the marked bound use $U_e$ with
$e=\lfloor d/(16b_0)\rfloor$. Its near-zero positivity covers
$[0,b_0/d]$. For sufficiently large $M$, $e\ge M^{1/5}$, and the
marked comparison applies with $e\asymp d$. Its expectation is
$O(1/d^2)$, by (4), proving the other assertion in \eqref{interval}.

At any row with $x\le a_0/(2d)$, all rows in the interval
in \eqref{interval} lie strictly before that row. Monotonicity and
the density relation give
\[
 (\ell+2)u_i(q)\frac c d
 \le\mu_1(a_0/d<x<b_0/d)\le C/d^2,
 \qquad M u_i(q)\le C/d.
\]
For $x\le a_0/(4M^{1/4})$ choose an integer $d$ comparable to
$\min(D,a_0/(2x))$, with the second argument infinite at zero.
It follows that
\[
 M u_i\le C(x+D^{-1})
       \le C(x+M^{-1}+\delta/\sqrt M).
\]
For larger $x$, the previous bound \eqref{mark} suffices, since
$\delta\le M^{-1/3}=o(M^{-1/4})\le Cx$.
Finally $x\le s$ and $M\asymp m$. The deleted index was arbitrary,
so the asserted improvement holds for every column.
\end{proof}

\begin{proof}[Proof of Theorem~\ref{main}]
The reviewed relative-balance theorem supplies \eqref{hyp} with
$a=1/3$. Lemma~\ref{step} improves its offset to $m^{-5/6}$.
A second application improves it to $m^{-1}$, proving the theorem.
Only two applications with fixed constants are involved.
\end{proof}
\end{document}
